% 2013 湖北卷理科第 12 题：程序框图（TikZ 矢量重绘）。
% 初值 a=10，i=1；判断 a=4？是则输出 i 并结束，否则按奇偶更新 a，再 i=i+1 返回。
\begin{tikzpicture}[font=\footnotesize, >=Stealth, line width=0.5pt,
  box/.style={draw, minimum height=1.9em, inner xsep=6pt},
  term/.style={draw, rounded corners=6pt, minimum height=1.9em, inner xsep=8pt},
  io/.style={draw, trapezium, trapezium left angle=70, trapezium right angle=110,
             minimum height=1.9em, inner xsep=6pt},
  dec/.style={draw, diamond, aspect=2.8, inner sep=1.5pt}]

  \node[term] (start) at (0,0) {开始};
  \node[box]  (init)  at (0,-1.2) {$a=10$，$i=1$};
  \node[dec]  (test)  at (0,-3.2) {$a=4$?};
  \node[dec]  (odd)   at (0,-6.4) {$a$ 是奇数?};
  \node[box]  (tri)   at (-2.4,-9.2) {$a=3a+1$};
  \node[box]  (half)  at (2.4,-9.2) {$a=\dfrac a2$};
  \node[box]  (inc)   at (0,-11.2) {$i=i+1$};
  \node[io]   (out)   at (5.4,-3.2) {输出 $i$};
  \node[term] (end)   at (5.4,-5.0) {结束};

  \draw[->] (start) -- (init);
  \draw[->] (init) -- (test);
  \draw[->] (test) -- node[left] {否} (odd);
  \draw[->] (test) -- node[above] {是} (out);
  \draw[->] (out) -- (end);
  \draw[->] (odd.west) -| node[pos=0.25, above] {是} (tri.north);
  \draw[->] (odd.east) -| node[pos=0.25, above] {否} (half.north);
  \draw[->] (tri.south) |- (inc.west);
  \draw[->] (half.south) |- (inc.east);
  \draw[->] (inc.south) -- (0,-12.4) -- (-4.6,-12.4) -- (-4.6,-2.2) -- (0,-2.2);

\end{tikzpicture}
